% ============================================================
%  课程笔记封面模板 —— 「墨竹图」风格
%  取材：文同《墨竹图》（北宋·湖州竹派）：素纸上以浓淡墨写数竿修竹，
%        竹节分明，撇叶如「个」「介」字，浓淡相破，不施粉黛。
%  与 cover_Impression风.tex 同构（眉题 / 主标题 / 菱形饰线 /
%  副标题 / 底部作者日期），直接 \input{} 复用：
%      % ============================================================
%  课程笔记封面模板 —— 「墨竹图」风格
%  取材：文同《墨竹图》（北宋·湖州竹派）：素纸上以浓淡墨写数竿修竹，
%        竹节分明，撇叶如「个」「介」字，浓淡相破，不施粉黛。
%  与 cover_Impression风.tex 同构（眉题 / 主标题 / 菱形饰线 /
%  副标题 / 底部作者日期），直接 \input{} 复用：
%      % ============================================================
%  课程笔记封面模板 —— 「墨竹图」风格
%  取材：文同《墨竹图》（北宋·湖州竹派）：素纸上以浓淡墨写数竿修竹，
%        竹节分明，撇叶如「个」「介」字，浓淡相破，不施粉黛。
%  与 cover_Impression风.tex 同构（眉题 / 主标题 / 菱形饰线 /
%  副标题 / 底部作者日期），直接 \input{} 复用：
%      % ============================================================
%  课程笔记封面模板 —— 「墨竹图」风格
%  取材：文同《墨竹图》（北宋·湖州竹派）：素纸上以浓淡墨写数竿修竹，
%        竹节分明，撇叶如「个」「介」字，浓淡相破，不施粉黛。
%  与 cover_Impression风.tex 同构（眉题 / 主标题 / 菱形饰线 /
%  副标题 / 底部作者日期），直接 \input{} 复用：
%      \input{cover_墨竹图风}
%
%  配色：mz* 前缀（素纸 + 墨分五色）。本文件用 \providecolor 兜底，
%        单独复制到任何笔记本也能编译；想整体换调改这里即可。
%
%  注意：整幅画面绘于 \begin{scope}[shift={(current page.south west)}] 内，
%        所有坐标以「页面左下角为原点、单位 cm」书写；脱离该 scope
%        直接写 (x,y)，图形会落到页面外而不可见。
%
%  可用字段（在主文件导言区用 \newcommand 预定义即可覆盖缺省值）：
%      \notename     主标题（必填，主模板已定义）
%      \noteauthor   作者（必填，主模板已定义）
%      \notecourseen 眉题英文（缺省 Course Notes）
%      \notesubtitle 副标题（缺省「学习笔记」）
%      \noteextra    底部附加信息（缺省不显示）
% ============================================================

% ---- 《墨竹图》取色（素纸 + 墨） ----
\providecolor{mzPaper}{RGB}{243,239,227}    % 素纸
\providecolor{mzPaperDark}{RGB}{224,218,201}% 素纸：暗部
\providecolor{mzInk3}{RGB}{166,164,156}     % 淡墨
\providecolor{mzInk2}{RGB}{98,98,94}        % 中墨
\providecolor{mzInk}{RGB}{44,44,42}         % 浓墨
\providecolor{mzSeal}{RGB}{168,54,46}       % 朱红钤印

% ---- 竹叶绘制宏：#1=颜色 #2=(x,y) #3=旋转角 #4=缩放 ----
\newcommand{\mzleaf}[4][mzInk]{%
	\begin{scope}[shift={#2}, rotate=#3, scale=#4]
		\fill[#1] (0,0) .. controls (0.72,0.44) and (2.05,0.30) .. (2.95,-0.45)
			.. controls (1.85,-0.12) and (0.62,0.10) .. (0,0) -- cycle;
	\end{scope}
}

% ---- 缺省字段（主文件已定义的同名命令不会被覆盖） ----
\providecommand{\notecourseen}{Course Notes}
\providecommand{\notesubtitle}{学\ 习\ 笔\ 记}
\providecommand{\noteextra}{}

\begin{titlepage}
	\begin{tikzpicture}[remember picture, overlay]
		\begin{scope}[shift={(current page.south west)}]

		% ============================================================
		%  1. 素纸（微渐变 + 斑驳）
		% ============================================================
		\shade[top color=mzPaper, bottom color=mzPaperDark] (0,0) rectangle (21,29.7);
		\foreach \x/\y in {2.6/27.2, 7.8/25.0, 12.6/28.0, 17.4/26.2, 4.2/19.6,
		                   13.8/15.8, 19.2/12.4, 8.4/9.2, 16.0/5.6}{
			\fill[mzInk3, opacity=0.08] (\x,\y) ellipse[x radius=1.0, y radius=0.38];
		}

		% ============================================================
		%  2. 背景淡墨渲染（衬托竹影）
		% ============================================================
		\fill[mzInk3, opacity=0.16]
			(1.2,10.6) .. controls (5.0,13.6) and (9.4,9.8) .. (13.6,12.6)
			.. controls (17.0,15.0) and (20.0,12.4) .. (20.6,9.4)
			-- (20.6,4.2) -- (1.2,4.2) -- cycle;

		% ============================================================
		%  3. 细竿（淡墨，居右）
		% ============================================================
		\draw[mzInk3, line width=3.0pt, line cap=round]
			(15.6,-0.6) .. controls (15.9,9.0) and (16.1,19.0) .. (16.4,30.4);
		\foreach \t in {0.14,0.32,0.50,0.68,0.86}{
			\draw[mzInk2, opacity=0.55, line width=1.5pt, line cap=round]
				({15.6+0.8*\t-0.28},{-0.6+31*\t}) -- ({15.6+0.8*\t+0.28},{-0.6+31*\t});
		}

		% ============================================================
		%  4. 中竿（中墨，居中）
		% ============================================================
		\draw[mzInk2, line width=4.6pt, line cap=round]
			(10.0,-0.6) .. controls (10.5,9.0) and (10.7,20.0) .. (11.2,30.4);
		\foreach \t in {0.10,0.26,0.42,0.58,0.74,0.90}{
			\draw[mzInk, opacity=0.7, line width=2.0pt, line cap=round]
				({10.0+1.2*\t-0.36},{-0.6+31*\t}) -- ({10.0+1.2*\t+0.36},{-0.6+31*\t});
		}

		% ============================================================
		%  5. 主竿（浓墨，居左，最粗）
		% ============================================================
		\draw[mzInk, line width=6.4pt, line cap=round]
			(5.2,-0.6) .. controls (5.7,9.0) and (6.0,20.0) .. (6.7,30.4);
		\foreach \t in {0.08,0.22,0.36,0.50,0.64,0.78,0.92}{
			\draw[mzInk, line width=2.6pt, line cap=round]
				({5.2+1.5*\t-0.48},{-0.6+31*\t}) -- ({5.2+1.5*\t+0.48},{-0.6+31*\t});
		}

		% ============================================================
		%  6. 撇叶（浓淡相破，如「个」「介」字）
		% ============================================================
		% 主竿：浓墨大叶
		\foreach \x/\y/\r/\s in {
				6.3/25.6/-28/1.20, 6.3/25.6/6/1.35,  6.3/25.6/40/1.05,
				6.0/19.8/-42/1.10, 6.0/19.8/-8/1.28, 6.0/19.8/28/0.98,
				5.7/13.4/-20/1.05, 5.7/13.4/12/1.20, 5.7/13.4/46/0.88}{
			\mzleaf{(\x,\y)}{\r}{\s}
		}
		% 中竿：中墨
		\foreach \x/\y/\r/\s in {
				10.6/23.2/-34/1.00, 10.6/23.2/2/1.15, 10.6/23.2/34/0.92,
				10.3/16.6/-24/0.95, 10.3/16.6/14/1.08, 10.3/16.6/44/0.85,
				10.1/10.2/-30/0.88, 10.1/10.2/8/1.00}{
			\mzleaf[mzInk2]{(\x,\y)}{\r}{\s}
		}
		% 细竿：淡墨小叶
		\foreach \x/\y/\r/\s in {
				16.1/21.4/-26/0.80, 16.1/21.4/6/0.92,  16.1/21.4/36/0.72,
				15.9/14.8/-18/0.76, 15.9/14.8/16/0.86}{
			\mzleaf[mzInk3]{(\x,\y)}{\r}{\s}
		}
		% 新篁（低矮的嫩竹两笔）
		\draw[mzInk2, line width=2.4pt, line cap=round]
			(18.4,-0.4) .. controls (18.6,2.2) and (18.8,4.4) .. (19.0,6.6);
		\mzleaf[mzInk2]{(19.0,6.6)}{-16}{0.72}
		\mzleaf[mzInk2]{(19.0,6.6)}{18}{0.80}

		% ============================================================
		%  7. 标题衬底（一层薄纸，保证文字可读）
		% ============================================================
		\fill[mzPaper, opacity=0.55]
			(2.6,16.9) -- (17.4,17.3) -- (18.0,22.6) -- (16.4,23.6) -- (2.2,22.9) -- cycle;

		% ============================================================
		%  8. 朱红钤印
		% ============================================================
		\fill[mzSeal, opacity=0.85] (18.3,25.0) rectangle (19.7,26.4);
		\draw[mzPaper, line width=0.7pt, opacity=0.9] (18.55,25.25) rectangle (19.45,26.15);
		\fill[mzSeal, opacity=0.78] (1.6,2.6) rectangle (2.8,3.8);
		\draw[mzPaper, line width=0.6pt, opacity=0.9] (1.82,2.82) rectangle (2.58,3.58);

		% ============================================================
		%  9. 细边框（墨线）
		% ============================================================
		\draw[mzInk, line width=0.7pt, opacity=0.42]
			([shift={(0.85,0.85)}]current page.south west) rectangle
			([shift={(-0.85,-0.85)}]current page.north east);

		% ============================================================
		% 10. 顶部眉题（课程英文小字 + 自适应墨线）
		% ============================================================
		\node[anchor=north, text=mzInk, align=center]
			(mzEyebrow) at ([yshift=-3.4cm]current page.north) {%
			{\sffamily\fontsize{12}{15}\selectfont \uppercase{\notecourseen}}
		};
		\draw[mzInk, line width=1.0pt] ($(mzEyebrow.west)+(-0.25cm,0)$) -- ++(-1.1cm,0);
		\draw[mzInk, line width=1.0pt] ($(mzEyebrow.east)+(0.25cm,0)$) -- ++(1.1cm,0);

		% ============================================================
		% 11. 主标题（居中，使用 \notename）
		% ============================================================
		\node[anchor=center, align=center, text=mzInk]
			at ([yshift=-8.3cm]current page.north) {%
			{\fontsize{38}{48}\sffamily\bfseries \notename}
		};

		% ============================================================
		% 12. 菱形 + 双细线
		% ============================================================
		\coordinate (C) at ([yshift=-11.6cm]current page.north);
		\draw[mzInk, line width=1.1pt] (C) ++(-2.2cm,0) -- ++(1.55cm,0);
		\draw[mzInk, line width=1.1pt] (C) ++(0.65cm,0) -- ++(1.55cm,0);
		\node[fill=mzInk, minimum size=0.24cm, inner sep=0pt, rotate=45] at (C) {};

		% ============================================================
		% 13. 副标题（使用 \notesubtitle）
		% ============================================================
		\node[anchor=north, text=mzInk, align=center]
			at ([yshift=-12.4cm]current page.north) {
			{\fontsize{15}{20}\sffamily \notesubtitle}
		};

		% ============================================================
		% 14. 底部作者、日期、附加信息
		% ============================================================
		\coordinate (B) at ([yshift=2.7cm]current page.south);
		\draw[mzInk, line width=0.8pt, opacity=0.8] (B) ++(-1.1cm,0.95cm) -- ++(2.2cm,0);
		\node[anchor=south, text=mzInk, align=center] at (B) {
			\begin{tabular}{c}
				{\fontsize{19}{24}\sffamily\bfseries \noteauthor} \\[0.35cm]
				{\large \sffamily \today}
				\ifstrempty{\noteextra}{}{\\[0.3cm]{\normalsize\sffamily \noteextra}}
			\end{tabular}
		};

		% ============================================================
		% 15. 右下角落款（画作信息）
		% ============================================================
		\node[anchor=south east, text=mzInk, opacity=0.6, font=\footnotesize\itshape]
			at ([shift={(-1.35,1.25)}]current page.south east)
			{Wen Tong, Ink Bamboo, c.1070s};

		\end{scope}
	\end{tikzpicture}
\end{titlepage}

%
%  配色：mz* 前缀（素纸 + 墨分五色）。本文件用 \providecolor 兜底，
%        单独复制到任何笔记本也能编译；想整体换调改这里即可。
%
%  注意：整幅画面绘于 \begin{scope}[shift={(current page.south west)}] 内，
%        所有坐标以「页面左下角为原点、单位 cm」书写；脱离该 scope
%        直接写 (x,y)，图形会落到页面外而不可见。
%
%  可用字段（在主文件导言区用 \newcommand 预定义即可覆盖缺省值）：
%      \notename     主标题（必填，主模板已定义）
%      \noteauthor   作者（必填，主模板已定义）
%      \notecourseen 眉题英文（缺省 Course Notes）
%      \notesubtitle 副标题（缺省「学习笔记」）
%      \noteextra    底部附加信息（缺省不显示）
% ============================================================

% ---- 《墨竹图》取色（素纸 + 墨） ----
\providecolor{mzPaper}{RGB}{243,239,227}    % 素纸
\providecolor{mzPaperDark}{RGB}{224,218,201}% 素纸：暗部
\providecolor{mzInk3}{RGB}{166,164,156}     % 淡墨
\providecolor{mzInk2}{RGB}{98,98,94}        % 中墨
\providecolor{mzInk}{RGB}{44,44,42}         % 浓墨
\providecolor{mzSeal}{RGB}{168,54,46}       % 朱红钤印

% ---- 竹叶绘制宏：#1=颜色 #2=(x,y) #3=旋转角 #4=缩放 ----
\newcommand{\mzleaf}[4][mzInk]{%
	\begin{scope}[shift={#2}, rotate=#3, scale=#4]
		\fill[#1] (0,0) .. controls (0.72,0.44) and (2.05,0.30) .. (2.95,-0.45)
			.. controls (1.85,-0.12) and (0.62,0.10) .. (0,0) -- cycle;
	\end{scope}
}

% ---- 缺省字段（主文件已定义的同名命令不会被覆盖） ----
\providecommand{\notecourseen}{Course Notes}
\providecommand{\notesubtitle}{学\ 习\ 笔\ 记}
\providecommand{\noteextra}{}

\begin{titlepage}
	\begin{tikzpicture}[remember picture, overlay]
		\begin{scope}[shift={(current page.south west)}]

		% ============================================================
		%  1. 素纸（微渐变 + 斑驳）
		% ============================================================
		\shade[top color=mzPaper, bottom color=mzPaperDark] (0,0) rectangle (21,29.7);
		\foreach \x/\y in {2.6/27.2, 7.8/25.0, 12.6/28.0, 17.4/26.2, 4.2/19.6,
		                   13.8/15.8, 19.2/12.4, 8.4/9.2, 16.0/5.6}{
			\fill[mzInk3, opacity=0.08] (\x,\y) ellipse[x radius=1.0, y radius=0.38];
		}

		% ============================================================
		%  2. 背景淡墨渲染（衬托竹影）
		% ============================================================
		\fill[mzInk3, opacity=0.16]
			(1.2,10.6) .. controls (5.0,13.6) and (9.4,9.8) .. (13.6,12.6)
			.. controls (17.0,15.0) and (20.0,12.4) .. (20.6,9.4)
			-- (20.6,4.2) -- (1.2,4.2) -- cycle;

		% ============================================================
		%  3. 细竿（淡墨，居右）
		% ============================================================
		\draw[mzInk3, line width=3.0pt, line cap=round]
			(15.6,-0.6) .. controls (15.9,9.0) and (16.1,19.0) .. (16.4,30.4);
		\foreach \t in {0.14,0.32,0.50,0.68,0.86}{
			\draw[mzInk2, opacity=0.55, line width=1.5pt, line cap=round]
				({15.6+0.8*\t-0.28},{-0.6+31*\t}) -- ({15.6+0.8*\t+0.28},{-0.6+31*\t});
		}

		% ============================================================
		%  4. 中竿（中墨，居中）
		% ============================================================
		\draw[mzInk2, line width=4.6pt, line cap=round]
			(10.0,-0.6) .. controls (10.5,9.0) and (10.7,20.0) .. (11.2,30.4);
		\foreach \t in {0.10,0.26,0.42,0.58,0.74,0.90}{
			\draw[mzInk, opacity=0.7, line width=2.0pt, line cap=round]
				({10.0+1.2*\t-0.36},{-0.6+31*\t}) -- ({10.0+1.2*\t+0.36},{-0.6+31*\t});
		}

		% ============================================================
		%  5. 主竿（浓墨，居左，最粗）
		% ============================================================
		\draw[mzInk, line width=6.4pt, line cap=round]
			(5.2,-0.6) .. controls (5.7,9.0) and (6.0,20.0) .. (6.7,30.4);
		\foreach \t in {0.08,0.22,0.36,0.50,0.64,0.78,0.92}{
			\draw[mzInk, line width=2.6pt, line cap=round]
				({5.2+1.5*\t-0.48},{-0.6+31*\t}) -- ({5.2+1.5*\t+0.48},{-0.6+31*\t});
		}

		% ============================================================
		%  6. 撇叶（浓淡相破，如「个」「介」字）
		% ============================================================
		% 主竿：浓墨大叶
		\foreach \x/\y/\r/\s in {
				6.3/25.6/-28/1.20, 6.3/25.6/6/1.35,  6.3/25.6/40/1.05,
				6.0/19.8/-42/1.10, 6.0/19.8/-8/1.28, 6.0/19.8/28/0.98,
				5.7/13.4/-20/1.05, 5.7/13.4/12/1.20, 5.7/13.4/46/0.88}{
			\mzleaf{(\x,\y)}{\r}{\s}
		}
		% 中竿：中墨
		\foreach \x/\y/\r/\s in {
				10.6/23.2/-34/1.00, 10.6/23.2/2/1.15, 10.6/23.2/34/0.92,
				10.3/16.6/-24/0.95, 10.3/16.6/14/1.08, 10.3/16.6/44/0.85,
				10.1/10.2/-30/0.88, 10.1/10.2/8/1.00}{
			\mzleaf[mzInk2]{(\x,\y)}{\r}{\s}
		}
		% 细竿：淡墨小叶
		\foreach \x/\y/\r/\s in {
				16.1/21.4/-26/0.80, 16.1/21.4/6/0.92,  16.1/21.4/36/0.72,
				15.9/14.8/-18/0.76, 15.9/14.8/16/0.86}{
			\mzleaf[mzInk3]{(\x,\y)}{\r}{\s}
		}
		% 新篁（低矮的嫩竹两笔）
		\draw[mzInk2, line width=2.4pt, line cap=round]
			(18.4,-0.4) .. controls (18.6,2.2) and (18.8,4.4) .. (19.0,6.6);
		\mzleaf[mzInk2]{(19.0,6.6)}{-16}{0.72}
		\mzleaf[mzInk2]{(19.0,6.6)}{18}{0.80}

		% ============================================================
		%  7. 标题衬底（一层薄纸，保证文字可读）
		% ============================================================
		\fill[mzPaper, opacity=0.55]
			(2.6,16.9) -- (17.4,17.3) -- (18.0,22.6) -- (16.4,23.6) -- (2.2,22.9) -- cycle;

		% ============================================================
		%  8. 朱红钤印
		% ============================================================
		\fill[mzSeal, opacity=0.85] (18.3,25.0) rectangle (19.7,26.4);
		\draw[mzPaper, line width=0.7pt, opacity=0.9] (18.55,25.25) rectangle (19.45,26.15);
		\fill[mzSeal, opacity=0.78] (1.6,2.6) rectangle (2.8,3.8);
		\draw[mzPaper, line width=0.6pt, opacity=0.9] (1.82,2.82) rectangle (2.58,3.58);

		% ============================================================
		%  9. 细边框（墨线）
		% ============================================================
		\draw[mzInk, line width=0.7pt, opacity=0.42]
			([shift={(0.85,0.85)}]current page.south west) rectangle
			([shift={(-0.85,-0.85)}]current page.north east);

		% ============================================================
		% 10. 顶部眉题（课程英文小字 + 自适应墨线）
		% ============================================================
		\node[anchor=north, text=mzInk, align=center]
			(mzEyebrow) at ([yshift=-3.4cm]current page.north) {%
			{\sffamily\fontsize{12}{15}\selectfont \uppercase{\notecourseen}}
		};
		\draw[mzInk, line width=1.0pt] ($(mzEyebrow.west)+(-0.25cm,0)$) -- ++(-1.1cm,0);
		\draw[mzInk, line width=1.0pt] ($(mzEyebrow.east)+(0.25cm,0)$) -- ++(1.1cm,0);

		% ============================================================
		% 11. 主标题（居中，使用 \notename）
		% ============================================================
		\node[anchor=center, align=center, text=mzInk]
			at ([yshift=-8.3cm]current page.north) {%
			{\fontsize{38}{48}\sffamily\bfseries \notename}
		};

		% ============================================================
		% 12. 菱形 + 双细线
		% ============================================================
		\coordinate (C) at ([yshift=-11.6cm]current page.north);
		\draw[mzInk, line width=1.1pt] (C) ++(-2.2cm,0) -- ++(1.55cm,0);
		\draw[mzInk, line width=1.1pt] (C) ++(0.65cm,0) -- ++(1.55cm,0);
		\node[fill=mzInk, minimum size=0.24cm, inner sep=0pt, rotate=45] at (C) {};

		% ============================================================
		% 13. 副标题（使用 \notesubtitle）
		% ============================================================
		\node[anchor=north, text=mzInk, align=center]
			at ([yshift=-12.4cm]current page.north) {
			{\fontsize{15}{20}\sffamily \notesubtitle}
		};

		% ============================================================
		% 14. 底部作者、日期、附加信息
		% ============================================================
		\coordinate (B) at ([yshift=2.7cm]current page.south);
		\draw[mzInk, line width=0.8pt, opacity=0.8] (B) ++(-1.1cm,0.95cm) -- ++(2.2cm,0);
		\node[anchor=south, text=mzInk, align=center] at (B) {
			\begin{tabular}{c}
				{\fontsize{19}{24}\sffamily\bfseries \noteauthor} \\[0.35cm]
				{\large \sffamily \today}
				\ifstrempty{\noteextra}{}{\\[0.3cm]{\normalsize\sffamily \noteextra}}
			\end{tabular}
		};

		% ============================================================
		% 15. 右下角落款（画作信息）
		% ============================================================
		\node[anchor=south east, text=mzInk, opacity=0.6, font=\footnotesize\itshape]
			at ([shift={(-1.35,1.25)}]current page.south east)
			{Wen Tong, Ink Bamboo, c.1070s};

		\end{scope}
	\end{tikzpicture}
\end{titlepage}

%
%  配色：mz* 前缀（素纸 + 墨分五色）。本文件用 \providecolor 兜底，
%        单独复制到任何笔记本也能编译；想整体换调改这里即可。
%
%  注意：整幅画面绘于 \begin{scope}[shift={(current page.south west)}] 内，
%        所有坐标以「页面左下角为原点、单位 cm」书写；脱离该 scope
%        直接写 (x,y)，图形会落到页面外而不可见。
%
%  可用字段（在主文件导言区用 \newcommand 预定义即可覆盖缺省值）：
%      \notename     主标题（必填，主模板已定义）
%      \noteauthor   作者（必填，主模板已定义）
%      \notecourseen 眉题英文（缺省 Course Notes）
%      \notesubtitle 副标题（缺省「学习笔记」）
%      \noteextra    底部附加信息（缺省不显示）
% ============================================================

% ---- 《墨竹图》取色（素纸 + 墨） ----
\providecolor{mzPaper}{RGB}{243,239,227}    % 素纸
\providecolor{mzPaperDark}{RGB}{224,218,201}% 素纸：暗部
\providecolor{mzInk3}{RGB}{166,164,156}     % 淡墨
\providecolor{mzInk2}{RGB}{98,98,94}        % 中墨
\providecolor{mzInk}{RGB}{44,44,42}         % 浓墨
\providecolor{mzSeal}{RGB}{168,54,46}       % 朱红钤印

% ---- 竹叶绘制宏：#1=颜色 #2=(x,y) #3=旋转角 #4=缩放 ----
\newcommand{\mzleaf}[4][mzInk]{%
	\begin{scope}[shift={#2}, rotate=#3, scale=#4]
		\fill[#1] (0,0) .. controls (0.72,0.44) and (2.05,0.30) .. (2.95,-0.45)
			.. controls (1.85,-0.12) and (0.62,0.10) .. (0,0) -- cycle;
	\end{scope}
}

% ---- 缺省字段（主文件已定义的同名命令不会被覆盖） ----
\providecommand{\notecourseen}{Course Notes}
\providecommand{\notesubtitle}{学\ 习\ 笔\ 记}
\providecommand{\noteextra}{}

\begin{titlepage}
	\begin{tikzpicture}[remember picture, overlay]
		\begin{scope}[shift={(current page.south west)}]

		% ============================================================
		%  1. 素纸（微渐变 + 斑驳）
		% ============================================================
		\shade[top color=mzPaper, bottom color=mzPaperDark] (0,0) rectangle (21,29.7);
		\foreach \x/\y in {2.6/27.2, 7.8/25.0, 12.6/28.0, 17.4/26.2, 4.2/19.6,
		                   13.8/15.8, 19.2/12.4, 8.4/9.2, 16.0/5.6}{
			\fill[mzInk3, opacity=0.08] (\x,\y) ellipse[x radius=1.0, y radius=0.38];
		}

		% ============================================================
		%  2. 背景淡墨渲染（衬托竹影）
		% ============================================================
		\fill[mzInk3, opacity=0.16]
			(1.2,10.6) .. controls (5.0,13.6) and (9.4,9.8) .. (13.6,12.6)
			.. controls (17.0,15.0) and (20.0,12.4) .. (20.6,9.4)
			-- (20.6,4.2) -- (1.2,4.2) -- cycle;

		% ============================================================
		%  3. 细竿（淡墨，居右）
		% ============================================================
		\draw[mzInk3, line width=3.0pt, line cap=round]
			(15.6,-0.6) .. controls (15.9,9.0) and (16.1,19.0) .. (16.4,30.4);
		\foreach \t in {0.14,0.32,0.50,0.68,0.86}{
			\draw[mzInk2, opacity=0.55, line width=1.5pt, line cap=round]
				({15.6+0.8*\t-0.28},{-0.6+31*\t}) -- ({15.6+0.8*\t+0.28},{-0.6+31*\t});
		}

		% ============================================================
		%  4. 中竿（中墨，居中）
		% ============================================================
		\draw[mzInk2, line width=4.6pt, line cap=round]
			(10.0,-0.6) .. controls (10.5,9.0) and (10.7,20.0) .. (11.2,30.4);
		\foreach \t in {0.10,0.26,0.42,0.58,0.74,0.90}{
			\draw[mzInk, opacity=0.7, line width=2.0pt, line cap=round]
				({10.0+1.2*\t-0.36},{-0.6+31*\t}) -- ({10.0+1.2*\t+0.36},{-0.6+31*\t});
		}

		% ============================================================
		%  5. 主竿（浓墨，居左，最粗）
		% ============================================================
		\draw[mzInk, line width=6.4pt, line cap=round]
			(5.2,-0.6) .. controls (5.7,9.0) and (6.0,20.0) .. (6.7,30.4);
		\foreach \t in {0.08,0.22,0.36,0.50,0.64,0.78,0.92}{
			\draw[mzInk, line width=2.6pt, line cap=round]
				({5.2+1.5*\t-0.48},{-0.6+31*\t}) -- ({5.2+1.5*\t+0.48},{-0.6+31*\t});
		}

		% ============================================================
		%  6. 撇叶（浓淡相破，如「个」「介」字）
		% ============================================================
		% 主竿：浓墨大叶
		\foreach \x/\y/\r/\s in {
				6.3/25.6/-28/1.20, 6.3/25.6/6/1.35,  6.3/25.6/40/1.05,
				6.0/19.8/-42/1.10, 6.0/19.8/-8/1.28, 6.0/19.8/28/0.98,
				5.7/13.4/-20/1.05, 5.7/13.4/12/1.20, 5.7/13.4/46/0.88}{
			\mzleaf{(\x,\y)}{\r}{\s}
		}
		% 中竿：中墨
		\foreach \x/\y/\r/\s in {
				10.6/23.2/-34/1.00, 10.6/23.2/2/1.15, 10.6/23.2/34/0.92,
				10.3/16.6/-24/0.95, 10.3/16.6/14/1.08, 10.3/16.6/44/0.85,
				10.1/10.2/-30/0.88, 10.1/10.2/8/1.00}{
			\mzleaf[mzInk2]{(\x,\y)}{\r}{\s}
		}
		% 细竿：淡墨小叶
		\foreach \x/\y/\r/\s in {
				16.1/21.4/-26/0.80, 16.1/21.4/6/0.92,  16.1/21.4/36/0.72,
				15.9/14.8/-18/0.76, 15.9/14.8/16/0.86}{
			\mzleaf[mzInk3]{(\x,\y)}{\r}{\s}
		}
		% 新篁（低矮的嫩竹两笔）
		\draw[mzInk2, line width=2.4pt, line cap=round]
			(18.4,-0.4) .. controls (18.6,2.2) and (18.8,4.4) .. (19.0,6.6);
		\mzleaf[mzInk2]{(19.0,6.6)}{-16}{0.72}
		\mzleaf[mzInk2]{(19.0,6.6)}{18}{0.80}

		% ============================================================
		%  7. 标题衬底（一层薄纸，保证文字可读）
		% ============================================================
		\fill[mzPaper, opacity=0.55]
			(2.6,16.9) -- (17.4,17.3) -- (18.0,22.6) -- (16.4,23.6) -- (2.2,22.9) -- cycle;

		% ============================================================
		%  8. 朱红钤印
		% ============================================================
		\fill[mzSeal, opacity=0.85] (18.3,25.0) rectangle (19.7,26.4);
		\draw[mzPaper, line width=0.7pt, opacity=0.9] (18.55,25.25) rectangle (19.45,26.15);
		\fill[mzSeal, opacity=0.78] (1.6,2.6) rectangle (2.8,3.8);
		\draw[mzPaper, line width=0.6pt, opacity=0.9] (1.82,2.82) rectangle (2.58,3.58);

		% ============================================================
		%  9. 细边框（墨线）
		% ============================================================
		\draw[mzInk, line width=0.7pt, opacity=0.42]
			([shift={(0.85,0.85)}]current page.south west) rectangle
			([shift={(-0.85,-0.85)}]current page.north east);

		% ============================================================
		% 10. 顶部眉题（课程英文小字 + 自适应墨线）
		% ============================================================
		\node[anchor=north, text=mzInk, align=center]
			(mzEyebrow) at ([yshift=-3.4cm]current page.north) {%
			{\sffamily\fontsize{12}{15}\selectfont \uppercase{\notecourseen}}
		};
		\draw[mzInk, line width=1.0pt] ($(mzEyebrow.west)+(-0.25cm,0)$) -- ++(-1.1cm,0);
		\draw[mzInk, line width=1.0pt] ($(mzEyebrow.east)+(0.25cm,0)$) -- ++(1.1cm,0);

		% ============================================================
		% 11. 主标题（居中，使用 \notename）
		% ============================================================
		\node[anchor=center, align=center, text=mzInk]
			at ([yshift=-8.3cm]current page.north) {%
			{\fontsize{38}{48}\sffamily\bfseries \notename}
		};

		% ============================================================
		% 12. 菱形 + 双细线
		% ============================================================
		\coordinate (C) at ([yshift=-11.6cm]current page.north);
		\draw[mzInk, line width=1.1pt] (C) ++(-2.2cm,0) -- ++(1.55cm,0);
		\draw[mzInk, line width=1.1pt] (C) ++(0.65cm,0) -- ++(1.55cm,0);
		\node[fill=mzInk, minimum size=0.24cm, inner sep=0pt, rotate=45] at (C) {};

		% ============================================================
		% 13. 副标题（使用 \notesubtitle）
		% ============================================================
		\node[anchor=north, text=mzInk, align=center]
			at ([yshift=-12.4cm]current page.north) {
			{\fontsize{15}{20}\sffamily \notesubtitle}
		};

		% ============================================================
		% 14. 底部作者、日期、附加信息
		% ============================================================
		\coordinate (B) at ([yshift=2.7cm]current page.south);
		\draw[mzInk, line width=0.8pt, opacity=0.8] (B) ++(-1.1cm,0.95cm) -- ++(2.2cm,0);
		\node[anchor=south, text=mzInk, align=center] at (B) {
			\begin{tabular}{c}
				{\fontsize{19}{24}\sffamily\bfseries \noteauthor} \\[0.35cm]
				{\large \sffamily \today}
				\ifstrempty{\noteextra}{}{\\[0.3cm]{\normalsize\sffamily \noteextra}}
			\end{tabular}
		};

		% ============================================================
		% 15. 右下角落款（画作信息）
		% ============================================================
		\node[anchor=south east, text=mzInk, opacity=0.6, font=\footnotesize\itshape]
			at ([shift={(-1.35,1.25)}]current page.south east)
			{Wen Tong, Ink Bamboo, c.1070s};

		\end{scope}
	\end{tikzpicture}
\end{titlepage}

%
%  配色：mz* 前缀（素纸 + 墨分五色）。本文件用 \providecolor 兜底，
%        单独复制到任何笔记本也能编译；想整体换调改这里即可。
%
%  注意：整幅画面绘于 \begin{scope}[shift={(current page.south west)}] 内，
%        所有坐标以「页面左下角为原点、单位 cm」书写；脱离该 scope
%        直接写 (x,y)，图形会落到页面外而不可见。
%
%  可用字段（在主文件导言区用 \newcommand 预定义即可覆盖缺省值）：
%      \notename     主标题（必填，主模板已定义）
%      \noteauthor   作者（必填，主模板已定义）
%      \notecourseen 眉题英文（缺省 Course Notes）
%      \notesubtitle 副标题（缺省「学习笔记」）
%      \noteextra    底部附加信息（缺省不显示）
% ============================================================

% ---- 《墨竹图》取色（素纸 + 墨） ----
\providecolor{mzPaper}{RGB}{243,239,227}    % 素纸
\providecolor{mzPaperDark}{RGB}{224,218,201}% 素纸：暗部
\providecolor{mzInk3}{RGB}{166,164,156}     % 淡墨
\providecolor{mzInk2}{RGB}{98,98,94}        % 中墨
\providecolor{mzInk}{RGB}{44,44,42}         % 浓墨
\providecolor{mzSeal}{RGB}{168,54,46}       % 朱红钤印

% ---- 竹叶绘制宏：#1=颜色 #2=(x,y) #3=旋转角 #4=缩放 ----
\newcommand{\mzleaf}[4][mzInk]{%
	\begin{scope}[shift={#2}, rotate=#3, scale=#4]
		\fill[#1] (0,0) .. controls (0.72,0.44) and (2.05,0.30) .. (2.95,-0.45)
			.. controls (1.85,-0.12) and (0.62,0.10) .. (0,0) -- cycle;
	\end{scope}
}

% ---- 缺省字段（主文件已定义的同名命令不会被覆盖） ----
\providecommand{\notecourseen}{Course Notes}
\providecommand{\notesubtitle}{学\ 习\ 笔\ 记}
\providecommand{\noteextra}{}

\begin{titlepage}
	\begin{tikzpicture}[remember picture, overlay]
		\begin{scope}[shift={(current page.south west)}]

		% ============================================================
		%  1. 素纸（微渐变 + 斑驳）
		% ============================================================
		\shade[top color=mzPaper, bottom color=mzPaperDark] (0,0) rectangle (21,29.7);
		\foreach \x/\y in {2.6/27.2, 7.8/25.0, 12.6/28.0, 17.4/26.2, 4.2/19.6,
		                   13.8/15.8, 19.2/12.4, 8.4/9.2, 16.0/5.6}{
			\fill[mzInk3, opacity=0.08] (\x,\y) ellipse[x radius=1.0, y radius=0.38];
		}

		% ============================================================
		%  2. 背景淡墨渲染（衬托竹影）
		% ============================================================
		\fill[mzInk3, opacity=0.16]
			(1.2,10.6) .. controls (5.0,13.6) and (9.4,9.8) .. (13.6,12.6)
			.. controls (17.0,15.0) and (20.0,12.4) .. (20.6,9.4)
			-- (20.6,4.2) -- (1.2,4.2) -- cycle;

		% ============================================================
		%  3. 细竿（淡墨，居右）
		% ============================================================
		\draw[mzInk3, line width=3.0pt, line cap=round]
			(15.6,-0.6) .. controls (15.9,9.0) and (16.1,19.0) .. (16.4,30.4);
		\foreach \t in {0.14,0.32,0.50,0.68,0.86}{
			\draw[mzInk2, opacity=0.55, line width=1.5pt, line cap=round]
				({15.6+0.8*\t-0.28},{-0.6+31*\t}) -- ({15.6+0.8*\t+0.28},{-0.6+31*\t});
		}

		% ============================================================
		%  4. 中竿（中墨，居中）
		% ============================================================
		\draw[mzInk2, line width=4.6pt, line cap=round]
			(10.0,-0.6) .. controls (10.5,9.0) and (10.7,20.0) .. (11.2,30.4);
		\foreach \t in {0.10,0.26,0.42,0.58,0.74,0.90}{
			\draw[mzInk, opacity=0.7, line width=2.0pt, line cap=round]
				({10.0+1.2*\t-0.36},{-0.6+31*\t}) -- ({10.0+1.2*\t+0.36},{-0.6+31*\t});
		}

		% ============================================================
		%  5. 主竿（浓墨，居左，最粗）
		% ============================================================
		\draw[mzInk, line width=6.4pt, line cap=round]
			(5.2,-0.6) .. controls (5.7,9.0) and (6.0,20.0) .. (6.7,30.4);
		\foreach \t in {0.08,0.22,0.36,0.50,0.64,0.78,0.92}{
			\draw[mzInk, line width=2.6pt, line cap=round]
				({5.2+1.5*\t-0.48},{-0.6+31*\t}) -- ({5.2+1.5*\t+0.48},{-0.6+31*\t});
		}

		% ============================================================
		%  6. 撇叶（浓淡相破，如「个」「介」字）
		% ============================================================
		% 主竿：浓墨大叶
		\foreach \x/\y/\r/\s in {
				6.3/25.6/-28/1.20, 6.3/25.6/6/1.35,  6.3/25.6/40/1.05,
				6.0/19.8/-42/1.10, 6.0/19.8/-8/1.28, 6.0/19.8/28/0.98,
				5.7/13.4/-20/1.05, 5.7/13.4/12/1.20, 5.7/13.4/46/0.88}{
			\mzleaf{(\x,\y)}{\r}{\s}
		}
		% 中竿：中墨
		\foreach \x/\y/\r/\s in {
				10.6/23.2/-34/1.00, 10.6/23.2/2/1.15, 10.6/23.2/34/0.92,
				10.3/16.6/-24/0.95, 10.3/16.6/14/1.08, 10.3/16.6/44/0.85,
				10.1/10.2/-30/0.88, 10.1/10.2/8/1.00}{
			\mzleaf[mzInk2]{(\x,\y)}{\r}{\s}
		}
		% 细竿：淡墨小叶
		\foreach \x/\y/\r/\s in {
				16.1/21.4/-26/0.80, 16.1/21.4/6/0.92,  16.1/21.4/36/0.72,
				15.9/14.8/-18/0.76, 15.9/14.8/16/0.86}{
			\mzleaf[mzInk3]{(\x,\y)}{\r}{\s}
		}
		% 新篁（低矮的嫩竹两笔）
		\draw[mzInk2, line width=2.4pt, line cap=round]
			(18.4,-0.4) .. controls (18.6,2.2) and (18.8,4.4) .. (19.0,6.6);
		\mzleaf[mzInk2]{(19.0,6.6)}{-16}{0.72}
		\mzleaf[mzInk2]{(19.0,6.6)}{18}{0.80}

		% ============================================================
		%  7. 标题衬底（一层薄纸，保证文字可读）
		% ============================================================
		\fill[mzPaper, opacity=0.55]
			(2.6,16.9) -- (17.4,17.3) -- (18.0,22.6) -- (16.4,23.6) -- (2.2,22.9) -- cycle;

		% ============================================================
		%  8. 朱红钤印
		% ============================================================
		\fill[mzSeal, opacity=0.85] (18.3,25.0) rectangle (19.7,26.4);
		\draw[mzPaper, line width=0.7pt, opacity=0.9] (18.55,25.25) rectangle (19.45,26.15);
		\fill[mzSeal, opacity=0.78] (1.6,2.6) rectangle (2.8,3.8);
		\draw[mzPaper, line width=0.6pt, opacity=0.9] (1.82,2.82) rectangle (2.58,3.58);

		% ============================================================
		%  9. 细边框（墨线）
		% ============================================================
		\draw[mzInk, line width=0.7pt, opacity=0.42]
			([shift={(0.85,0.85)}]current page.south west) rectangle
			([shift={(-0.85,-0.85)}]current page.north east);

		% ============================================================
		% 10. 顶部眉题（课程英文小字 + 自适应墨线）
		% ============================================================
		\node[anchor=north, text=mzInk, align=center]
			(mzEyebrow) at ([yshift=-3.4cm]current page.north) {%
			{\sffamily\fontsize{12}{15}\selectfont \uppercase{\notecourseen}}
		};
		\draw[mzInk, line width=1.0pt] ($(mzEyebrow.west)+(-0.25cm,0)$) -- ++(-1.1cm,0);
		\draw[mzInk, line width=1.0pt] ($(mzEyebrow.east)+(0.25cm,0)$) -- ++(1.1cm,0);

		% ============================================================
		% 11. 主标题（居中，使用 \notename）
		% ============================================================
		\node[anchor=center, align=center, text=mzInk]
			at ([yshift=-8.3cm]current page.north) {%
			{\fontsize{38}{48}\sffamily\bfseries \notename}
		};

		% ============================================================
		% 12. 菱形 + 双细线
		% ============================================================
		\coordinate (C) at ([yshift=-11.6cm]current page.north);
		\draw[mzInk, line width=1.1pt] (C) ++(-2.2cm,0) -- ++(1.55cm,0);
		\draw[mzInk, line width=1.1pt] (C) ++(0.65cm,0) -- ++(1.55cm,0);
		\node[fill=mzInk, minimum size=0.24cm, inner sep=0pt, rotate=45] at (C) {};

		% ============================================================
		% 13. 副标题（使用 \notesubtitle）
		% ============================================================
		\node[anchor=north, text=mzInk, align=center]
			at ([yshift=-12.4cm]current page.north) {
			{\fontsize{15}{20}\sffamily \notesubtitle}
		};

		% ============================================================
		% 14. 底部作者、日期、附加信息
		% ============================================================
		\coordinate (B) at ([yshift=2.7cm]current page.south);
		\draw[mzInk, line width=0.8pt, opacity=0.8] (B) ++(-1.1cm,0.95cm) -- ++(2.2cm,0);
		\node[anchor=south, text=mzInk, align=center] at (B) {
			\begin{tabular}{c}
				{\fontsize{19}{24}\sffamily\bfseries \noteauthor} \\[0.35cm]
				{\large \sffamily \today}
				\ifstrempty{\noteextra}{}{\\[0.3cm]{\normalsize\sffamily \noteextra}}
			\end{tabular}
		};

		% ============================================================
		% 15. 右下角落款（画作信息）
		% ============================================================
		\node[anchor=south east, text=mzInk, opacity=0.6, font=\footnotesize\itshape]
			at ([shift={(-1.35,1.25)}]current page.south east)
			{Wen Tong, Ink Bamboo, c.1070s};

		\end{scope}
	\end{tikzpicture}
\end{titlepage}
